% Compile with: xelatex ceos_clos_wiring.tex   (or lualatex)
% Requires a Unicode monospace font with box-drawing glyph coverage.
% DejaVu Sans Mono (used below) is free and covers U+2500-257F fully.
\documentclass[11pt]{article}
\usepackage{fontspec}
\setmainfont{DejaVu Sans}
\setmonofont{DejaVu Sans Mono}
\usepackage[margin=1.2cm,landscape,a3paper]{geometry}
\usepackage{booktabs}
\usepackage{array}
\usepackage{titlesec}
\titleformat{\section}{\normalfont\Large\bfseries}{\thesection}{1em}{}

\begin{document}

\section*{EVPN-VXLAN Lab --- ceos\_clos.yml}
\subsection*{Layer 1 wiring, loopbacks and interfaces}

{\small
\textit{Note: the ten spine--leaf uplinks are shown as labeled point-to-point stubs at each}\\
\textit{node rather than as ten crossing diagonal lines, to keep the diagram readable. Each}\\
\textit{stub names both the local port and the exact remote node:port it terminates on.}
}

\begin{verbatim}
┌──────────────┐          ┌──────────────┐
│    spine1    │          │    spine2    │
│ 10.0.0.21/32 │          │ 10.0.0.22/32 │
└──────────────┘          └──────────────┘

eth1 -> leaf1:eth1        eth1 -> leaf1:eth2
eth2 -> leaf2:eth1        eth2 -> leaf2:eth2
eth3 -> leaf3:eth1        eth3 -> leaf3:eth2
eth4 -> leaf4:eth1        eth4 -> leaf4:eth2
eth5 -> leaf5:eth1        eth5 -> leaf5:eth2


eth1->spine1            eth1->spine1            eth1->spine1            eth1->spine1            eth1->spine1
eth2->spine2            eth2->spine2            eth2->spine2            eth2->spine2            eth2->spine2
┌──────────────┐        ┌──────────────┐        ┌──────────────┐        ┌──────────────┐        ┌──────────────┐
│    leaf1     │        │    leaf2     │        │    leaf3     │        │    leaf4     │        │    leaf5     │
│ 10.0.0.11/32 │        │ 10.0.0.12/32 │        │ 10.0.0.13/32 │        │ 10.0.0.14/32 │        │ 10.0.0.15/32 │
└───────┬──────┘        └───────┬──────┘        └───────┬──────┘        └───────┬──────┘        └───────┬──────┘
        │ eth3                  │ eth3                  │ eth3                  │ eth3                  │ eth3
        │                       │                       │                       │                       │
        │                       │                       │                       │                       │
        │                       │                       │                       │                       │
        │                       │                       │                       │                       │
        └──────┐        ┌───────┘                       │                       │                       │
          eth1 │        │ eth2                          │ eth1                  │ eth1                  │ eth1
            ┌──┴────────┴──┐                    ┌───────┴──────┐        ┌───────┴──────┐        ┌───────┴──────┐
            │    host1     │                    │    host2     │        │    host3     │        │    host4     │
            │192.168.10.31 │                    │192.168.10.32 │        │192.168.10.33 │        │192.168.20.34 │
            └──────────────┘                    └──────────────┘        └──────────────┘        └──────────────┘


# --- OVERLAY SUBNETS ---
  host1, host2, host3  ->  VLAN 10 / L2VNI 10010  ->  192.168.10.0/24  (VRF TENANT_A)
  host4                ->  VLAN 20 / L2VNI 10020  ->  192.168.20.0/24  (VRF TENANT_B)

# --- HOST ATTACHMENTS ---
    host1: Dual-homed to leaf1 + leaf2 -> ESI-LAG / EVPN multihoming lab
    host2 / host3: Single-homed, same VLAN/VNI on different leaves -> L2 EVPN stretch lab (Type-2 routes)
    host4: Single-homed, different VLAN/VNI/VRF -> L3 inter-VNI EVPN routing lab (IRB / Type-5)

\end{verbatim}

\newpage
\subsection*{Loopback0 addresses}

\begin{tabular}{ll}
\toprule
\textbf{Node} & \textbf{Loopback0} \\
\midrule
spine1 & 10.0.0.21/32 \\
spine2 & 10.0.0.22/32 \\
leaf1  & 10.0.0.11/32 \\
leaf2  & 10.0.0.12/32 \\
leaf3  & 10.0.0.13/32 \\
leaf4  & 10.0.0.14/32 \\
leaf5  & 10.0.0.15/32 \\
\bottomrule
\end{tabular}

\vspace{0.4cm}

\noindent
\begin{tabular}{lll}
\toprule
\textbf{Host} & \textbf{Attached leaf(s)} & \textbf{Data IP} \\
\midrule
host1 & leaf1 + leaf2 (ESI-LAG bond0) & 192.168.10.31/24 \\
host2 & leaf3                         & 192.168.10.32/24 \\
host3 & leaf4                         & 192.168.10.33/24 \\
host4 & leaf5                         & 192.168.20.34/24 \\
\bottomrule
\end{tabular}

\subsection*{Overlay (host-facing) subnets}

\begin{tabular}{llllll}
\toprule
\textbf{VLAN / VNI} & \textbf{Subnet} & \textbf{Anycast GW} & \textbf{VRF} & \textbf{Leaves} & \textbf{L3VNI} \\
\midrule
VLAN10 / L2VNI 10010 & 192.168.10.0/24 & 192.168.10.1 & TENANT\_A & leaf1--4 & 50010 \\
VLAN20 / L2VNI 10020 & 192.168.20.0/24 & 192.168.20.1 & TENANT\_B & leaf5    & 50020 \\
\bottomrule
\end{tabular}

\vspace{0.4cm}

\noindent
\begin{tabular}{lll}
\toprule
\textbf{Host} & \textbf{Attached leaf(s)} & \textbf{Data IP} \\
\midrule
host1 & leaf1 + leaf2 (ESI-LAG bond0) & 192.168.10.31/24 \\
host2 & leaf3                         & 192.168.10.32/24 \\
host3 & leaf4                         & 192.168.10.33/24 \\
host4 & leaf5                         & 192.168.20.34/24 \\
\bottomrule
\end{tabular}

\subsection*{RD / RT scheme}

\begin{tabular}{lll}
\toprule
\textbf{Scope} & \textbf{RD pattern} & \textbf{RT (common, must match to import)} \\
\midrule
VLAN10 (per leaf) & \texttt{<leaf-loopback>:10}  & 10:10 \\
VLAN20 (leaf5)     & \texttt{<leaf-loopback>:20}  & 20:20 \\
VRF TENANT\_A       & \texttt{<leaf-loopback>:100} & 100:100 \\
VRF TENANT\_B       & \texttt{<leaf-loopback>:200} & 200:200 \\
\bottomrule
\end{tabular}

\vspace{0.2cm}

{\small
\noindent
\textit{Note: TENANT\_A and TENANT\_B use different RTs on purpose --- no route leaking, so host4 will not reach host1/2/3.}\\
\textit{That's intended VRF isolation, not a bug --- good talking point for the interview.}
}

\subsection*{ESI (ethernet-segment) for host1's dual-homed link}

\begin{tabular}{lll}
\toprule
\textbf{Leaf} & \textbf{Port} & \textbf{ESI} \\
\midrule
leaf1 & Port-Channel1 & \texttt{0000:0000:0000:1001:0001} \\
leaf2 & Port-Channel1 & \texttt{0000:0000:0000:1001:0001} (same ESI = same Ethernet Segment) \\
\bottomrule
\end{tabular}

\end{document}
